\clearpage\thispagestyle{empty}\null\clearpage
\TFMChapter[even]{Discusión}
\TFMSection{Lecciones Aprendidas durante Desarrollo de la Solución}

\TFMSection{Riesgos, Ética y Mitigaciones}
El carácter de prueba de concepto del proyecto no exime de analizar sus riesgos como si fuera a manejar información y clientes reales, ya que muchas de las decisiones tomadas tendrían que revisarse precisamente antes de dar ese paso. Esta sección recoge los riesgos identificados durante el desarrollo, agrupados en tres bloques, junto con las mitigaciones ya incorporadas y las que quedan pendientes.

\subsection{Seguridad}
El sistema separa las credenciales de los usuarios finales de las credenciales de servicio que utiliza internamente el pipeline documental para persistir metadatos, de forma que un empleado nunca comparte su sesión con un proceso automático; las contraseñas se almacenan cifradas y el acceso a la API requiere un token de sesión con caducidad. No obstante, se han identificado dos simplificaciones de seguridad propias de una fase de prueba de concepto que deberían corregirse antes de cualquier despliegue con datos reales: la política de acceso entre orígenes (CORS) de la API está completamente abierta, y el manejador de errores expone el detalle técnico completo de la excepción cuando el sistema se ejecuta en modo de depuración, lo que podría filtrar información interna si ese modo quedase activo en producción por error. Asimismo, la infraestructura en AWS se ha configurado manualmente en la consola en lugar de mediante una definición de infraestructura como código, lo que dificulta auditar y reproducir de forma fiable los permisos concedidos a cada componente.

\subsection{Ética y protección de datos}
El sistema procesa datos personales del cliente final (nombre, fechas de viaje, preferencias) y documentación comercial de los proveedores turísticos, y los envía a servicios de terceros (proveedores de modelos de lenguaje) para su procesamiento. Esto implica una cesión de datos a terceros que, en un despliegue real y no solo como prueba de concepto, debería quedar cubierta por una base legal explícita y por los acuerdos de tratamiento de datos correspondientes con cada proveedor. Por otro lado, al tratarse de un sistema que genera contenido de forma automática (contexto turístico, itinerario, propuesta comercial), existe un riesgo inherente de que el modelo introduzca información incorrecta o inventada. Este riesgo se mitiga en el diseño mediante dos mecanismos: la fundamentación de las respuestas en documentación real a través del sistema RAG, con un umbral mínimo de relevancia antes de dar por válido un fragmento, y la revisión humana obligatoria de la propuesta final antes de su publicación, de forma que ninguna propuesta llega al cliente sin que una persona la haya validado.

\subsection{Riesgos operativos y técnicos}
Durante la exploración del código se han identificado varios puntos que, sin afectar al funcionamiento observado del sistema, conviene documentar como riesgo técnico: el módulo de reordenación semántica de resultados (\textit{reranking}) del sistema RAG está implementado pero no se invoca en la ruta de consulta activa, por lo que la calidad de la recuperación podría mejorarse sin necesidad de desarrollo adicional; el estado de una conversación interrumpida en un punto de revisión humana se conserva en memoria del proceso del orquestador, por lo que un reinicio del contenedor en ese momento concreto podría perder el progreso de esa conversación (aunque las propuestas ya completadas sí quedan a salvo en la base de datos); y existen dos agentes (de redacción y de maquetación) implementados en una iteración anterior del diseño que ya no forman parte del flujo activo del grafo, lo que representa código muerto a eliminar o documentar explícitamente. Todos estos puntos se recogen como líneas de trabajo futuro en el Capítulo 8.
